\part{The state of the programme}
\chapter{Solved, open, impossible}\label{chap:state}

This chapter consolidates the closing remarks of the chapters. Verdicts are quoted, not
re-argued; the chapter and statement numbers point to where the argument is.

Four of the open problems of the first edition of this list have since been settled, three
of them affirmatively and one by a no-go theorem, and a fifth has been settled in a way that
was not on the menu at all. Composite $q$ is proved (Chapter~\ref{chap:XVII}); the rank-$d$
unitarity mechanism is proved, and then made unconditional (Chapters~\ref{chap:XVIII},
\ref{chap:XX}); the search for a cyclic class computing the $L$-invariant is closed by an
obstruction (Chapter~\ref{chap:XIX}); and the bidirected sandpile theory asked for by the
assembly chapters turns out not to exist, for a reason that is itself a theorem
(Chapter~\ref{chap:Bio7}).

\section{Corrections made by later chapters to earlier ones}

\begin{itemize}
\item The untwisted boundary type is $\mathrm{III}_{1/q}$ outright, not a
$\beta$-family: Chapter~\ref{chap:III}, Remark~1.6 corrects Chapter~\ref{chap:I},
Theorem~2.4(iii). The applied note of Chapter~\ref{chap:Holo1} sharpens this further:
$\mathrm{III}_{1/q^2}$ when the graph is bipartite.
\item The line-digraph sandpile recursion is $|\Jac(Y_{k+1})|=q^{\,n_k(q-1)-1}|\Jac(Y_k)|$:
Chapter~\ref{chap:III}, Theorem~2.1 corrects the exponent $2^{n_k-1}$ of
Chapter~\ref{chap:II}, Remark~5.2 (the two agree for $q=2$). This recursion should be
compared with Levine's theorem on sandpile groups of line digraphs (J.\ Combin.\ Theory
Ser.\ A 118 (2011)); the series flags the comparison and does not claim priority.
\item The naive $\Z_p[[T]]$ conjecture for the congruence tower
(Chapter~\ref{chap:II}, Conjecture~3.6(ii-b)) is withdrawn in the same chapter and in
Chapter~\ref{chap:III}; three independent certificates against it are $\lambda=-1$
(Ch.~\ref{chap:III}), $A_kA_k^{\mathsf T}\ne A_k^{\mathsf T}A_k$ (Ch.~\ref{chap:V}), and
the torsion-freeness of the $p$-part of the pro-module with $\mu>0$
(Ch.~\ref{chap:VIII}).
\item The shadow of a $q$-regular level needs $q-1$ return arcs, not $q$:
Chapter~\ref{chap:VIII}, Remark~1.4.
\item The geodesic edge flow must be defined by ``not in a common chamber'', not ``not
adjacent'': Chapter~\ref{chap:XV}, Remark~1.2 corrects Chapter~\ref{chap:XI},
Definition~3.1; the two agree only on flag complexes, and of the five test complexes
only $Y_{168}$ is a flag complex.
\item The $p$-primary hypothesis in the tail-norm theorem was an artefact of the proof, not
of the statement: Chapter~\ref{chap:XVII}, Theorem~\ref{XVII:thm:main} removes it for every
$q\ge2$, because the rank of a line-digraph adjacency matrix is $n$ over \emph{every} field
(Ch.~\ref{chap:XVII}, Lemma~\ref{XVII:lem:rank}), which is what the counting argument needs.
\item The invariant subspace left open at Chapter~\ref{chap:XVIII}, Remark~4.2(a) does
exist: Chapter~\ref{chap:XX}, Corollary~\ref{XX:cor:W}. The candidate intertwiner recorded
there was the right one, and the identity it needs is a single flag count.
\item The expectation, stated in the assembly chapters, that double-strandedness would be
handled by a sandpile theory on a quotient graph is withdrawn:
Chapter~\ref{chap:Bio7}, Theorem~\ref{Bio7:thm:anti}. Chapter~\ref{chap:Bio6} had already
shown that no new lemma is \emph{needed}, by passing to the double cover instead; Chapter
\ref{chap:Bio7} shows that the quotient route is not merely unnecessary but unavailable.
\item Cross-references between the original papers drifted as sections were
renumbered (for instance ``[V, Prob.~5.1]'' is Problem~4.1 of Chapter~\ref{chap:V};
``[I, Prop.~1.7(iii)]'' is Proposition~1.9(iii) of Chapter~\ref{chap:I}). The
concordance and the chapter-wise numbering of this book are authoritative.
\item The type III framing of ``holographic failure'' and the claim that the control
positive-realization problem inherits the trace obstruction were both wrong and are
retracted in the applied chapters themselves (Chapters~\ref{chap:Holo1},
\ref{chap:Ctrl1}).
\end{itemize}

\section{Proved}

\begin{itemize}
\item The boundary channel is arithmetically blind: $K_0(C(\partial X)\rtimes\Gamma)
\cong\Z^r\oplus\Z/(r-1)$ (Ch.~\ref{chap:I}, Thm.~3.7).
\item The trace obstruction: no non-negative matrix, of any size, is shift equivalent to
the Ihara companion (Ch.~\ref{chap:II}, Thms.~1.2, 1.5).
\item The shadow realization: $K_0(\mathcal O_{B(Y)})\cong\Z\oplus\Jac(Y)$, $K_1\cong\Z$,
functorially and at every level of the tower (Ch.~\ref{chap:VI}, Thm.~2.1;
Ch.~\ref{chap:VIII}, Thm.~1.3); the signed-shadow lemma for any integer matrix with
positive diagonal (Ch.~\ref{chap:X}, Lem.~5.1).
\item The abelian limit module $\Lambda^n/D(1+T)$ with exact control
(Ch.~\ref{chap:IV}, Thm.~2.1); the exceptional zero of order two and the formula
$\tfrac12\Theta''(0)=-\kappa(Y_0)\|h_\alpha\|^2$ (Ch.~\ref{chap:IX}, Thm.~1.2); the five
faces of the $L$-invariant (Ch.~\ref{chap:XIV}, Cor.~7.2).
\item The degeneracy calculus and the Euler rank of the degeneracy complex
(Ch.~\ref{chap:V}, Prop.~1.2, Thm.~2.1); the tail norm and its pro-module
(Ch.~\ref{chap:V}, Thm.~3.1); vanishing of the upper Hecke operator on the norm kernel
and the snake factorization (Ch.~\ref{chap:VII}, Lem.~1.1, Thm.~2.1); the Eisenstein
boundary sequence $0\to\Z/n_k\to\coker\Delta_k^0\to\Jac(Y_k)\to0$ and its
non-splitting on $K_4$ (Ch.~\ref{chap:XIII}, Thm.~3.2).
\item \textbf{For every $q\ge2$}, $\ker(\tau_*|\Jac(Y_{k+1}))=\Jac(Y_{k+1})[q]\cong
(\Z/q)^{r_k}$, and every invariant factor divisible by a prime $p\mid q$ is divisible by
$p^{v_p(q)}$ (Ch.~\ref{chap:XVII}, Thm.~\ref{XVII:thm:main}); hence the pro-module has no
$q$-torsion (Ch.~\ref{chap:XVII}, Cor.~\ref{XVII:cor:notorsion}) and the sequence
$0\to\mathfrak X_\infty^\vee\to K_0(\mathcal O_\infty)\to\Z[1/q]\to0$ splits for every $q$
(Ch.~\ref{chap:XVII}, Thm.~\ref{XVII:thm:split}). The $p$-prime case is
Ch.~\ref{chap:VIII}, Thm.~2.8; the inductive limit is Kirchberg with
$\operatorname{Tor}K_0$ dual to the pro-module (Ch.~\ref{chap:VIII}, Thms.~3.2, 3.3); the
head correspondence with $K_0$-index $\eta_*$ and the Hecke round trips
(Ch.~\ref{chap:XVI}, Thms.~2.3, 3.1, 4.1).
\item In rank two: no cubic trace obstruction (Ch.~\ref{chap:X}, Thm.~4.2); $b_1=0$ for
thick quotients, hence no abelian towers (Ch.~\ref{chap:XII}, Thm.~1.1); the covering
norm theorem and the incidence calculus (Ch.~\ref{chap:XII}, Thms.~2.1, 3.1); the
edge-flow intertwiner $L_E\Psi=\Psi C'$ (Ch.~\ref{chap:XI}, Thm.~3.3); the two-sided
reduction $\ell^2(E)=\operatorname{im}\Psi\oplus W$ and the unitarity of the remainder,
$RR^{\mathsf T}=qI$ (Ch.~\ref{chap:XV}, Thms.~2.1, 2.2, 3.1).
\item \textbf{In every rank $d\ge2$}: the unitarity identity
$L_EL_E^{\mathsf T}=q^{\,d-1}I+q^{\,d-1}(q-1)T^{\mathsf T}T$ and its dual
(Ch.~\ref{chap:XVIII}, Thm.~\ref{XVIII:thm:main}), whose whole content is that two
hyperplanes of $\mathrm{PG}(d,q)$ always meet in codimension two; exactly $|E|-n$ singular
values equal $q^{(d-1)/2}$ and $L_E:\ker S\to\ker T$ is a bijective isometry
(Ch.~\ref{chap:XVIII}, Cor.~\ref{XVIII:cor:uncond}). The invariant subspace that makes the
critical circle unconditional is supplied in Ch.~\ref{chap:XX}: the number of complete flags
of $\mathrm{PG}(d,q)$ with $F_{i+1}=V$ and $F_1$ inside a hyperplane $H$ takes exactly two
values according as $V\subseteq H$ or not (Thm.~\ref{XX:thm:flag}), which \emph{is} the
intertwining $L_EN_i^{\mathsf T}=q^{\,i}(N_0^{\mathsf T}A_{i+1}-N_{i+1}^{\mathsf T})$
(Thm.~\ref{XX:thm:inter}); hence $W_d=(\operatorname{im}\Psi_d)^\perp$ is two-sided
invariant and every eigenvalue of $L_E|_{W_d}$ has modulus $q^{(d-1)/2}$ with no hypothesis
whatever (Cor.~\ref{XX:cor:W}). The companion matrix is explicit in every rank and its
characteristic polynomial is the local Hecke polynomial of $\mathrm{PGL}_{d+1}$
(Thm.~\ref{XX:thm:hecke}).
\item In the applications: the bundle-chain splitting theorem
$\K(G)\cong\bigoplus(\Z/r_c)^{\ell_c-1}\oplus\K(\mathrm{Sk})$ and the canonicity of the
skeleton (Ch.~\ref{chap:Bio3}); the two-family theorem $\K(D_w)=\Z/c$
(Ch.~\ref{chap:Bio4}); the power-sum bound on hidden eigenvalues
(Ch.~\ref{chap:Ctrl1}); the sign of $\lambda$ as a normality certificate
(Ch.~\ref{chap:Net1}).
\item In double-stranded assembly: the reverse-complement double cover is balanced, so the
whole single-strand apparatus applies to it verbatim and the five plus-strand and two
minus-strand ribosomal operons of \emph{E.\ coli} merge into one seven-copy family
(Ch.~\ref{chap:Bio6}, Thm.~\ref{Bio6:thm:merge}, Cor.~\ref{Bio6:cor:rrna}); the quotient is
a signed graph exactly when no repeated $k$-mer is its own reverse complement, which is
automatic for odd $k$ (Ch.~\ref{chap:Bio7}, Prop.~\ref{Bio7:prop:parity}); it is unbalanced
exactly when the sequence carries an inverted repeat of length at least $k$
(Thm.~\ref{Bio7:thm:balance}); and the Reiner--Tseng covering sequence holds on it, with a
non-split $2$-primary part arising from a real bacteriophage (Thm.~\ref{Bio7:thm:rt}).
\end{itemize}

\section{Verified computationally, not proved}

The parahoric factorization on $Y_{168}$ (Ch.~\ref{chap:XI}, Cor.~3.5) was numeric
before Ch.~\ref{chap:XV} proved it in general; the Ramanujan property of the five
$\widetilde A_2$ test complexes is numeric to $10^{-6}$ (Ch.~\ref{chap:X}, Thm.~4.5(i));
the character mechanism for kernel orders along finite $\widetilde A_2$ chains is
verified to $10^{-6}$ (Ch.~\ref{chap:XII}, Thm.~2.1(iii)); finiteness of
$K_0(\mathcal O_{L_E})$ rests on $1\notin\operatorname{spec}L_E$ on $Y_{168}$
(Ch.~\ref{chap:XI}, Thm.~3.7); the non-commutativity certificate
$A_kA_k^{\mathsf T}\ne A_k^{\mathsf T}A_k$ is checked on $K_4$ only (Ch.~\ref{chap:V},
Prop.~3.4); the exhaustive statements about chord diagrams hold for $m\le5$
(Ch.~\ref{chap:Bio3}, Rem.~9) and are known to fail at $m=6$.

Two items are new. The corank of $\Psi_d$ is at least $d(d+1)$ by
Ch.~\ref{chap:XX}, Theorem~\ref{XX:thm:corank}, and equals it on all three thick
$\widetilde A_2$ complexes tested; equality for $d\ge3$ is unproved, and with it the exact
value of $\dim W_d$. And a signed $r$-bundle chain of $s$ internal vertices raises the
$r$-rank by exactly $s$ in every case tested (Ch.~\ref{chap:Bio7},
Thm.~\ref{Bio7:thm:series}); only the inequality is proved.

\section{Open}

\begin{enumerate}
\item \textbf{The main conjecture for the degeneracy algebra.} Construct, from
$L$-values alone, an element $L_H\in K_1(H_\Sigma)$ equal to the characteristic class
$[\Theta_H]$ modulo the image of $K_1(H)$ (Ch.~\ref{chap:XIII}, Rem.~5.1); before that,
determine the image of $K_1(H)$, i.e.\ how non-unique the class is, and whether
$H\to H_\Sigma$ is stably flat (Ch.~\ref{chap:XIII}, Rem.~5.2). This is Problem~4.1 of
Ch.~\ref{chap:V} in its final form, and it is now the oldest unsettled question in the
programme.
\item \textbf{Rank two realization.} A non-negative integer matrix shift equivalent
over $\Z$ to the graded companion $C^{(3)}$ --- an instance of the weak Generalized
Spectral Conjecture of Boyle--Handelman (Ch.~\ref{chap:XI}, Rem.~2.4).
\item \textbf{Congruence towers in rank two.} The growth law of $\Jac^H$ along a genuine
infinite congruence tower of $\widetilde A_2$ complexes; no second level has been
computed (Ch.~\ref{chap:XII}, Rem.~6.1).
\item \textbf{The size of the unitary part.} Whether
$\operatorname{corank}\Psi_d=d(d+1)$ for $d\ge3$, and whether $W_d$ is the \emph{largest}
two-sided invariant subspace inside $\ker S\cap\ker T$ (Ch.~\ref{chap:XX},
Rem.~\ref{XX:rem:soi}). Also a thick $\widetilde A_3$ complex small enough to compute with,
and an exact certificate of the vertex Ramanujan property of the test complexes
(Ch.~\ref{chap:XV}, Rem.~6.1(c)).
\item \textbf{Counting double-stranded reconstructions.} The reverse-complement involution
induces a perfect \emph{symmetric} linking form on $\K(D_k)$
(Ch.~\ref{chap:Bio7}, Rem.~\ref{Bio7:rem:pairing}). At tree level it is a bijection from
the arborescences converging on a root to those diverging from the image root, not an
action, so there are no orbits to sum over; and for $\phi$X174 the group order is not a
square, so there is no metabolizer either. What a double-stranded count is built from is
open.
\item \textbf{Non-pairwise arrangement invariants.} Any presentation of the arrangement
group for repeats with three or more copies that is not built from pairwise data
(Ch.~\ref{chap:Bio4}, Open~(a)); and the even-$k$ case of the quotient, which is a signed
graph with half-edges branched over the palindromic repeated $k$-mers
(Ch.~\ref{chap:Bio7}, Rem.~\ref{Bio7:rem:halfedge}).
\end{enumerate}

\section{Impossible}

No non-negative matrix or dilation realizes the Ihara companion; no Hecke-dynamical
$C^*$-algebra carries $\Z\oplus\Jac(Y)$ as $K_0$ with the Ihara flow as gauge dynamics;
no separable $C^*$-algebra has $K_0\supseteq\mathfrak X_\infty$; no $\Z_p[[T]]$-module
produces $\lambda=-1$; no unital $*$-homomorphism $\mathcal O_{k+1}\to\mathcal O_k$
induces $\eta_*$; no cyclic cocycle on the commutative Hecke--Dirac triple computes the
$L$-invariant; no $L$-invariant exists for the Iwahori tower; no $\Z_p$-tower of thick
$\widetilde A_2$ complexes exists; no cubic trace obstruction exists; no invariant of the
unsigned chord-crossing graph computes the arrangement group of two-copy repeats, and no
pairwise invariant at all computes it for three-copy repeats; no coset structure
underlies long-read constraints in the reconstruction torsor.

Two entries are added here, and both were open problems on the previous list.

\emph{No odd cyclic class computes the graph $L$-invariant by a linear construction.}
The crossed product proposed for the job is Morita-trivial --- the deck action on a
connected $\Z$-cover is free, so $C_0(\widetilde Y_\alpha)\rtimes\Z\cong\bigoplus_V\mathcal
K$ and $K_1=0$ (Ch.~\ref{chap:XIX}, Prop.~\ref{XIX:prop:compacts}). In the Floquet
commutant, where $K_1$ is not zero, the natural element does not exist, because the pencil
degenerates exactly at the point one wants to evaluate at; what the odd pairing reads there
is the \emph{order} of the exceptional zero, which is two for every voltage class, whereas
$\kappa_0\|h_\alpha\|^2$ is its leading \emph{coefficient}
(Ch.~\ref{chap:XIX}, Prop.~\ref{XIX:prop:notinv}, Cor.~\ref{XIX:cor:winding}). And in any
algebra whose $K_1$ has rank one, a pairing linear in each slot is a product of two linear
forms and therefore vanishes on a subspace, while $\kappa_0\|h_\alpha\|^2$ is positive
definite; so no such construction can work once $b_1\ge2$
(Ch.~\ref{chap:XIX}, Thm.~\ref{XIX:thm:obstruction}). What is \emph{not} excluded is an
algebra with $\operatorname{rank}K_1\ge b_1$: the reduced group $C^*$-algebra of
$\pi_1(Y)$ is free of rank $b_1$ and sits exactly at that threshold.

\emph{The directed sandpile group of a reverse-complement de Bruijn graph has no covering
theory, and its signed quotient does not contract.} With $P$ the permutation matrix of the
involution, $P\Delta P=\Delta^{\mathsf T}$ and never $\Delta$: the deck involution is an
anti-automorphism of the Laplacian, so the $\pm1$ eigenlattices are not $\Delta$-stable and
the classical covering factorization is unavailable
(Ch.~\ref{chap:Bio7}, Thm.~\ref{Bio7:thm:anti}). The Reiner--Tseng sequence that does hold
computes an invariant of the undirected shadow, which by the matrix-tree theorem counts
spanning trees; the quantity of interest counts Eulerian circuits and lives on the digraph.
Separately, a signed $r$-bundle chain obeys the series law --- the presentation column at
the chain's endpoint differs from the contracted graph by $s\,r\,d_0$ --- so
$\K(\Sigma)\not\cong(\Z/r)^s\oplus\K(\Sigma/C)$, and the chain contraction that makes a
bacterial genome computable on one strand has no bidirected analogue
(Ch.~\ref{chap:Bio7}, Thm.~\ref{Bio7:thm:series}, Cor.~\ref{Bio7:cor:nocompute}).
